\documentclass[11pt]{article}
\usepackage{icmlko}
\begin{document}

\title{문턱은 일곱 언저리다 --- 그리고 깊이 열둘의 손해는 이미 갚혀 있었다}
\author{월드모델 실험실}
\date{노트 522 · 2026년 8월 4일}
\maketitle

\begin{abstract}
노트 521 이 ``깊이만 자료 쪽 갈래를 연다''를 남겼다. 어느 깊이에서 켜지는지
쟀다.

\begin{center}
\begin{tabular}{lrl}
\hline
깊이 & \texttt{TRAINR} 판 이득 & 부호 \\
\hline
$6$ & $\mathbf{-0.0000}$ & $5/12$ \\
$\mathbf{8}$ & $\mathbf{+0.0018}$ & $\mathbf{12/12}$ \\
$10$ & $+0.0024$ & $11/12$ \\
$12$ & $\mathbf{+0.0028}$ & $10/12$ \\
$14$ & $+0.0019$ & $11/12$ \\
\hline
\end{tabular}
\end{center}

\textbf{문턱이 날카롭다} --- $6$ 에서 정확히 $0$($5/12$, 동전 던지기)이고
$8$ 에서 $+0.0018$ 인데 부호가 $12/12$ 다. 그 뒤로는 평평한 봉우리($12$)를
지나 $14$ 에서 다시 준다.
\end{abstract}

\begin{figure}[t]
\centering
\includegraphics[width=\textwidth]{figs/f1.pdf}
\caption{\textbf{가} --- 깊이 곡선. \textbf{나} --- 깊이 $12$ 를 다시
세기. \textbf{다} --- 손해를 언제 적을 수 있나.}
\end{figure}

\section{깊이 열둘의 값을 다시 센다}

노트 457 이 깊이를 $6 \to 12$ 로 올리며 적었다 --- \emph{``판은 확실히
잃고(} $-0.0032$ \emph{) 집 밖 두 짝을 산다(} $+0.029$ \emph{). 환율
$1.95$. 거래를 숨기지 않는다.''}

\textbf{그 손해는 이미 갚혀 있었다.} 깊이 $12$ 가 연 문 뒤에
\texttt{TRAINR} $+0.0028$ 이 있었고, 육십다섯 노트 동안 아무도 안 가져갔다.

\begin{center}
\begin{tabular}{lr}
\hline
& 판 \\
\hline
노트 457 이 적은 깊이 $12$ 의 손해 & $-0.0032$ \\
그것이 연 문 뒤의 \texttt{TRAINR} & $\mathbf{+0.0028}$ \\
\hline
진짜 순값 & $\mathbf{-0.0004}$ \\
\hline
\end{tabular}
\end{center}

\textbf{깊이 $12$ 는 판을 거의 안 잃었고, 집 밖 $+0.029$ 는 사실상
공짜였다.} 그때 적은 환율 $1.95$ 는 \emph{그 시점에 아는 것으로는} 맞았지만
\emph{실제로는} 훨씬 좋은 거래였다.

\begin{quote}
\textbf{조항} --- \textbf{채택이 판을 잃으면, 그 채택이 \emph{여는} 갈래를
세어 보고 나서 손해를 적는다.} 용량을 올리는 채택은 그 용량에서 죽었던
자료 쪽 후보 목록을 함께 남기고, 그것을 다 재기 전에는 손해가 확정이 아니다.
\end{quote}

\section{문턱이 날카롭다는 것}

$6$ 과 $8$ 사이에서 $0 \to +0.0018$ 이고 부호가 $5/12 \to 12/12$ 로
바뀐다. 크기보다 \textbf{부호 일치가 먼저 켜진다} --- 깊이 $8$ 의 이득은
$12$ 의 $64\%$ 인데 부호는 더 고르다.

읽으면 --- 최근 가중이 사는 것은 \textbf{일곱 차 언저리의 상호작용}이다.
그보다 얕으면 어느 행을 얼마나 보든 같은 것만 표현할 수 있고, 그 위에서는
표현할 수 있는 것이 갑자기 늘어난다.

\textbf{정직하게} --- 깊이는 트리의 \emph{최대} 깊이지 실제 쓰인 차수가
아니다. 노트 478 이 잎 깊이 중앙값을 $11.0$ 으로 쟀다. ``일곱 차''는
상한의 이름이지 센 값이 아니다.

\section*{남은 것}

\begin{itemize}
\item \textbf{다른 자료 후보도 깊이 감응인가} --- \texttt{TRAINW} $\cdot$
\texttt{TRAINCAP} $\cdot$ \texttt{DOMDROP} 을 깊이 $6$ 에서 재면 같은
그림인가. 같다면 \textbf{이 사이클 다섯 채택이 전부 노트 457 의 자식}이다.
\item \textbf{깊이 $12$ 를 다시 고를까} --- 문턱 위 평평한 구간에서 $12$ 가
봉우리이고 $14$ 가 내려간다. 지금 값이 맞다.
\item \textbf{원장에 ``이 채택이 여는 갈래'' 칸} --- 노트 $519 \cdot 520$
이 두 번 제안했고 이 노트가 셋째다. \textbf{사람이 안 잊게 하는 유일한
방법이다.}
\end{itemize}

\end{document}
